\documentclass[12pt,a4paper]{article}

\usepackage[utf8]{inputenc}
\usepackage[T1]{fontenc}
\usepackage{lmodern}
\usepackage[margin=1in]{geometry}
\usepackage{setspace}
\usepackage[hidelinks]{hyperref}
\usepackage{amsmath,amssymb}
\usepackage{enumitem}
\usepackage{booktabs}
\usepackage{titlesec}
\usepackage[stretch=15]{microtype}
\usepackage{parskip}
\usepackage[backend=biber,style=authoryear,maxnames=3,minnames=1]{biblatex}
\addbibresource{references.bib}

% Section formatting
\titleformat{\section}{\large\bfseries}{}{0em}{}
\titleformat{\subsection}{\normalsize\bfseries}{}{0em}{}
\titleformat{\subsubsection}{\normalsize\itshape}{}{0em}{}

% Math shortcuts
\newcommand{\E}{\mathbb{E}}
\newcommand{\Var}{\mathrm{Var}}
\newcommand{\SD}{\mathrm{SD}}
\newcommand{\IG}{\mathrm{IG}}
\newcommand{\Normal}{\mathcal{N}}

\title{Technical Note: Literature Review of the\\Variance Decomposition
Framework\\[6pt]
\large Prior Art for $\rho$, $\kappa$, and the Law of Total Variance\\
as a Prior Specification Tool\\[6pt]
\normalsize Current state on branch \texttt{feat/joe-restructure}}
\author{}
\date{February 2026}

\begin{document}
\maketitle
\thispagestyle{empty}

\vspace{1em}
\hrule
\vspace{1em}

\section{Purpose}

The paper introduces a variance decomposition for hyperprior construction
(Section~3.1) built on three elements:

\begin{enumerate}[label=(\roman*)]
  \item The \textbf{law of total variance} used as a design tool to
    analytically derive hyperprior parameters from a target unconditional
    (direct) prior:
    \[
      v^* = \Var{m(\psi)} + \E{v(\psi)}
    \]
  \item The \textbf{variance partition ratio} $\rho \in (0,1)$, defined as the
    fraction of direct prior variance attributable to uncertainty about the
    population center:
    \[
      \rho\, v^* = \Var{m(\psi)}, \qquad (1-\rho)\, v^* = \E{v(\psi)}
    \]
  \item The \textbf{coefficient of variation of heterogeneity}
    $\kappa = \SD[\sigma^2] / \E{\sigma^2}$, a distribution-free measure of
    uncertainty about the level of population heterogeneity.
\end{enumerate}

This note surveys existing literature to determine whether any of these three
elements have been proposed before, and how to position the paper's
contribution relative to closely related work.


\section{The Variance Partition Ratio $\rho$}

\subsection{Post-hoc diagnostics: VPC and ICC}

The idea of expressing variance as a ratio of between-group to total variance
is well established --- as a \emph{posterior diagnostic}. Two standard terms
exist:

\begin{itemize}
  \item \textbf{Intraclass Correlation Coefficient (ICC).} Dating to Fisher
    (1925) and earlier, the ICC is defined as the correlation between two
    randomly chosen units within the same group. In a simple two-level
    random-intercept model, the ICC equals
    $\sigma^2_u / (\sigma^2_u + \sigma^2_e)$.

  \item \textbf{Variance Partition Coefficient (VPC).} Introduced by Goldstein,
    Browne \& Rasbash (2002) in \emph{Understanding Statistics}. The VPC
    generalises the ICC to settings where the ICC interpretation (as a
    correlation) breaks down: models with random slopes, heteroscedastic
    residuals, non-normal responses, and three or more levels. Goldstein et
    al.\ define the VPC as the proportion of total variance attributable to a
    given level, conditional on covariates.
\end{itemize}

In both cases the quantity is computed \emph{from a fitted model} to summarise
how much clustering exists in the data. The standard references are Snijders \&
Bosker (2012, \emph{Multilevel Analysis}, 2nd ed.) and Leckie et al.\ (2020)
for extensions to count data. The meta-analysis analogue is the $I^2$ statistic
of Higgins \& Thompson (2002), defined as
$\tau^2 / (\sigma^2 + \tau^2)$.

\medskip

\noindent\textbf{Relationship to $\rho$.}\quad The paper's $\rho$ is
numerically the same object as the VPC/ICC --- but used in a fundamentally
different way. In the existing literature, VPC/ICC is a \emph{descriptive
summary} of an estimated model. In the paper, $\rho$ is a \emph{design input}:
a researcher choice that, together with the direct prior moments $(m^*, v^*)$,
analytically determines the hyperprior parameters. No prior work uses VPC/ICC
as a specification device for constructing priors.

An important related observation: Natarajan \& Kass (2000, \emph{JASA}) showed
that the commonly used inverse-gamma$(\varepsilon, \varepsilon)$ prior on
individual variance components induces a bimodal prior on the ICC/VPC,
concentrated near 0 and 1 --- an extremely informative (and usually
unintended) implied prior on the variance partition. Specifying $\rho$ directly
avoids this pathology entirely, because the partition is a fixed researcher
choice rather than an emergent consequence of diffuse priors on individual
variance components.


\subsection{Variance proportions for prior specification}

A recent and growing literature does use variance proportions as inputs for
prior specification, rather than post-hoc diagnostics. This is the most
relevant body of work.

\subsubsection{Fuglstad, Hem, Knight, Rue \& Riebler (2020)}

``Intuitive Joint Priors for Variance Parameters,'' \emph{Bayesian Analysis}
15(4):1109--1137. This is the single closest existing paper.

Fuglstad et al.\ reparameterise a model's variance components into a
\emph{total variance} and a set of \emph{proportions} $(\omega_1, \ldots,
\omega_K)$ that allocate variance to different model components along a tree
structure. At each split in the tree, the analyst specifies either a Dirichlet
prior (expressing ignorance about the split) or a penalised complexity (PC)
prior (expressing a preference for one branch). The software implementation is
the \texttt{makemyprior} R package (Hem, Fuglstad \& Riebler, 2024,
\emph{Journal of Statistical Software}).

\medskip
\noindent\textbf{Three key differences from $\rho$:}

\begin{enumerate}[label=(\alph*)]
  \item \textbf{What is decomposed.} Fuglstad et al.\ decompose variance
    \emph{across model components} at the same level of the hierarchy (e.g.,
    spatial effect vs.\ temporal effect vs.\ residual). The paper's $\rho$
    decomposes variance \emph{across hierarchical levels} via the law of total
    variance identity $V = \Var[\E(\vartheta|\psi)] + \E[\Var(\vartheta|\psi)]$.
    These are different decompositions: one is additive across components, the
    other is the between/within split of the iterated expectations identity.

  \item \textbf{Random vs.\ fixed.} In the HD prior framework, the proportions
    $\omega_j$ are \emph{random variables} with their own priors (Dirichlet or
    PC). In the paper, $\rho$ is a \emph{fixed researcher choice} --- a scalar
    that, together with the direct prior moments, deterministically maps to
    hyperprior parameters. This is a deliberate design decision: $\rho$
    encodes a structural belief about the hierarchy, not a quantity to be
    learned from the data.

  \item \textbf{Starting point.} Fuglstad et al.\ specify priors directly on
    the reparameterised variance components. They do not begin with a target
    unconditional (direct) prior on an individual's quantity and solve backwards
    for hyperprior parameters. The paper's approach uses the direct prior as the
    anchor and derives everything else from it.
\end{enumerate}

\subsubsection{Aguilar \& B\"{u}rkner (2023)}

``Intuitive Joint Priors for Bayesian Linear Multilevel Models: The R2D2M2
Prior,'' \emph{Electronic Journal of Statistics} 17(1):1319--1374.

The R2D2M2 prior extends the R2D2 shrinkage prior (Zhang et al., 2022,
\emph{JASA}) to multilevel models. The user places a Beta prior on the model's
$R^2$ (coefficient of determination), then decomposes explained variance across
regression terms via a Dirichlet distribution. The hyperparameters are
$\mu_{R^2}$ (expected $R^2$) and $\phi_{R^2}$ (precision/confidence in that
expectation).

\medskip
\noindent\textbf{Relationship to the paper.}\quad The R2D2M2 framework shares
the broad goal of making variance-related hyperparameters interpretable. The
$\phi_{R^2}$ parameter captures ``how confident the user is about $R^2$,''
which is conceptually adjacent to $\kappa$ (how confident the researcher is
about the level of heterogeneity). However, R2D2M2 operates in the regression
setting (decomposing explained variance across predictors), not in the
hierarchical random-effects setting (decomposing unconditional variance across
hierarchical levels). The objects being reparameterised are different.

\subsubsection{Ferrari \& Ventrucci (2025)}

``A Standardisation Procedure to Incorporate Variance Partitioning Based Priors
in Latent Gaussian Models,'' arXiv:2501.16057. This very recent paper extends
the Fuglstad et al.\ framework to a broader class of latent Gaussian models,
confirming that variance-partition-based prior specification is an active
research area. The same distinctions (model components vs.\ hierarchical
levels, random proportions vs.\ fixed choices) apply.


\section{The Coefficient of Variation of Heterogeneity $\kappa$}

The specific construction
\[
  \kappa \coloneqq \frac{\SD[\sigma^2]}{\E{\sigma^2}}
\]
as a named, interpretable, distribution-free hyperparameter for specifying
uncertainty about the level of population heterogeneity does not appear in
the existing literature. This appears to be a genuinely novel contribution.

\subsection{What is known}

The mathematical fact that the inverse-gamma distribution $\IG(\alpha, \beta)$
has coefficient of variation $1/\sqrt{\alpha - 2}$ (for $\alpha > 2$) is
textbook material. It appears on the Wikipedia page for the inverse-gamma
distribution and in standard references. The relationship
$\kappa = 1/\sqrt{\alpha_0 - 2}$, equivalently $\alpha_0 = 2 + 1/\kappa^2$,
is therefore well known \emph{as a property of the IG distribution}.

What is new is:
\begin{enumerate}[label=(\roman*)]
  \item Defining $\kappa$ in a \textbf{distribution-free} way --- as
    $\SD[\sigma^2]/\E{\sigma^2}$ without reference to any specific family for
    $\sigma^2$ --- and then showing that it maps to $\alpha_0$ \emph{if} the IG
    family is chosen.
  \item Proposing $\kappa$ as a \textbf{primary specification device}: the
    researcher thinks in terms of $\kappa$ (``how uncertain am I about the
    level of heterogeneity?'') rather than $\alpha_0$ (``what are the
    degrees-of-freedom of my inverse-gamma prior?'').
  \item Providing an \textbf{interpretive scale}: $\kappa \to 0$ means
    certainty about heterogeneity, $\kappa = 1$ means the standard deviation of
    $\sigma^2$ equals its mean, and $\kappa > 1$ means very high uncertainty.
\end{enumerate}

\subsection{Closest parallels}

\begin{itemize}
  \item \textbf{$\phi_{R^2}$ in R2D2M2} (Aguilar \& B\"{u}rkner, 2023). A
    precision parameter for the Beta prior on $R^2$, capturing confidence in the
    expected $R^2$ value. Conceptually related (both capture ``meta-uncertainty''
    about a variance-related quantity), but $\phi_{R^2}$ is a Beta precision
    parameter in the regression setting, not a CV of a variance component in the
    hierarchical setting.

  \item \textbf{$\alpha_0$ as ``prior sample size.''} The standard textbook
    interpretation of the IG shape parameter is as an effective sample size.
    This is an interpretation of $\alpha_0$ itself, not a reparameterisation
    into a unit-free, distribution-free quantity.

  \item \textbf{The $I^2$ statistic} (Higgins \& Thompson, 2002). A
    distribution-free heterogeneity measure in meta-analysis. However, $I^2$
    measures the \emph{proportion} of total variance from heterogeneity
    (analogous to $\rho$), not the \emph{uncertainty about the level of
    heterogeneity} (which is what $\kappa$ captures). They are complementary
    concepts.
\end{itemize}


\section{The Law of Total Variance as a Prior Specification Tool}

The paper uses the law of total variance (equivalently, the law of iterated
expectations combined with the variance decomposition identity) as an
\emph{analytical tool} for deriving hyperprior parameters. Given a target
unconditional variance $v^*$, the identity
\[
  v^* = \underbrace{\Var[\E(\vartheta_i|\psi)]}_{\rho\, v^*}
      + \underbrace{\E[\Var(\vartheta_i|\psi)]}_{(1-\rho)\, v^*}
\]
is solved for the hyperprior parameters that produce the target moments. This
use of the identity --- as a design equation, not a descriptive one --- does
not appear in the existing literature.

\subsection{Closest existing work}

\subsubsection{Manderson \& Goudie (2023)}

``Translating Predictive Distributions into Informative Priors,''
arXiv:2303.08528. On first reading this appears to be the closest existing
work. Three successive drafts of this note have revised the comparison, each
time arriving at a sharper understanding. The current assessment follows.

\medskip
\noindent\textbf{Their setup.}\quad Manderson \& Goudie work with a
\emph{single-level} model: an observable $Y$, parameters $\theta$, and
hyperparameters $\lambda$ that parameterise the prior $P(\theta|\lambda)$.
Their goal is to find $\lambda$ such that the \emph{prior predictive
distribution}
\[
  P(Y|\lambda) = \int P(Y|\theta)\, P(\theta|\lambda)\, d\theta
\]
matches a target distribution $T(Y)$ elicited from a domain expert. There is
one integration: over parameters $\theta$, given fixed hyperparameters
$\lambda$. The quantity being matched is the distribution of \emph{observables}
$Y$, not of parameters. Their examples are all single-level: a cure fraction
survival model, an $R^2$ prior for linear regression, and a human growth curve
model. None involves a hierarchical population structure.

\medskip
\noindent\textbf{The correct alignment.}\quad The present paper has two
distinct steps, and Manderson \& Goudie correspond to the \emph{first} step,
not the second:

\begin{center}
\begin{tabular}{lll}
\toprule
\textbf{Step} & \textbf{This paper} & \textbf{Manderson \& Goudie} \\
\midrule
Beliefs about quantities & Jacobian correction & Distribution matching \\
$\to$ prior on $\theta$ & (analytical, \S3.1) & (numerical) \\
\addlinespace
Direct prior & Variance decomposition & \emph{Not addressed} \\
$\to$ BHM hyperpriors & with $(\rho, \kappa)$ (\S3.2) & \\
\bottomrule
\end{tabular}
\end{center}

Both the Jacobian correction and Manderson \& Goudie's distribution matching
solve the same problem: ``I have beliefs about derived/observable quantities;
translate these into a prior on model parameters.'' The Jacobian correction does
this analytically when the mapping $g: \theta \to \vartheta$ is injective and
differentiable. Manderson \& Goudie do it numerically for arbitrary models,
matching the full prior predictive distribution via Bayesian optimisation. They
are alternative solutions to the \emph{same} step.

The paper's \emph{second} step --- the variance decomposition --- has no
counterpart in Manderson \& Goudie. This is the step that takes a direct prior
on $\vartheta$ for a single individual and derives hyperprior parameters $\xi$
for a BHM's population distribution:
\[
  p^*(\vartheta) = \int p(\vartheta|\psi)\, p(\psi|\xi)\, d\psi
\]
Manderson \& Goudie do not address this integral at all. They have no
population structure, no between/within decomposition, and no analogue of
$\rho$ or $\kappa$.

\medskip
\noindent\textbf{The identifiability problem.}\quad Both approaches face
non-uniqueness, but at different steps and for different reasons.

Manderson \& Goudie face it at the quantity-to-prior step: they note that
``there often are many optimal values $\lambda^*$ that yield values of
$\tilde{D}(\lambda^*|\mathbf{X})$ that are practically
indistinguishable\ldots with immensely differing prior distributions.'' Their
resolution is generic regularisation (maximise prior variance) with
multi-objective optimisation over a Pareto frontier. The paper's Jacobian
correction avoids this entirely: given an injective $g$ and a density
$p_\vartheta$, the induced prior $p_\theta$ is unique.

The present paper faces the identifiability problem at the \emph{second} step:
the direct prior $p^*(\vartheta)$ is characterised by two moments $(m^*, v^*)$,
but the NIG hyperprior has four free parameters
$(\mu_0, \tau_0^2, \alpha_0, \beta_0)$. Two equations in four unknowns --- the
system is underdetermined. The paper's resolution is to supply two additional
\emph{interpretable} constraints:
\begin{itemize}
  \item $\rho$ pins down the split between $\tau_0^2$ and $\E{\sigma^2}$:
    how much of the direct prior variance comes from uncertainty about the
    population mean vs.\ genuine heterogeneity.
  \item $\kappa$ pins down $\SD[\sigma^2]/\E{\sigma^2}$: how confident the
    researcher is about the level of heterogeneity.
\end{itemize}
Together, $(m^*, v^*, \rho, \kappa)$ provide exactly four constraints for four
unknowns, and the solution is closed-form. This replaces regularisation with
\emph{substantive meaning}: each degree of freedom corresponds to an
interpretable belief about the population structure.

\medskip
\noindent\textbf{Composability.}\quad The two approaches operate at different
levels and could in principle be composed into a full pipeline:
\begin{enumerate}
  \item Use Manderson \& Goudie (or the Jacobian correction) to go from beliefs
    about observables/generated quantities to a direct prior on parameters.
  \item Use the variance decomposition with $(\rho, \kappa)$ to go from the
    direct prior to BHM hyperpriors.
\end{enumerate}
Neither paper alone covers both steps. Manderson \& Goudie stop after step~1;
the present paper assumes step~1 is done (via the Jacobian correction or by
direct specification) and provides step~2.

\subsubsection{Kadane et al.\ (1980) and Gabry et al.\ (2019)}

Kadane et al.\ (1980) established the principle that prior elicitation should
be conducted on \emph{observable quantities} rather than abstract parameters.
Gabry, Simpson, Vehtari, Betancourt \& Gelman (2019, ``Visualization in
Bayesian Workflow'') advocate for prior predictive checks as a central part of
Bayesian practice. Gelman's associated writing explicitly endorses the idea that
``you can include prior information on quantities of interest, not just on
parameters in your model.''

Both articulate the \emph{principle} but operationalise it only through
simulation and visual calibration. The paper provides the analytical complement:
exact formulas for what hyperpriors must be, given the target marginal.

\subsubsection{Gelman (2006) and Simpson et al.\ (2017)}

Gelman (2006, ``Prior Distributions for Variance Parameters in Hierarchical
Models,'' \emph{Bayesian Analysis}) addresses which \emph{family} of prior to
use on variance parameters (half-Cauchy, half-$t$, uniform on SD), not how to
determine the \emph{parameters} of a chosen family from a target unconditional
distribution.

Simpson, Rue, Riebler, Martins \& S{\o}rbye (2017, ``Penalising Model
Component Complexity,'' \emph{Statistical Science}) define PC priors via KL
divergence from a base model, with the user making a single tail probability
statement. This is a prior specification method, but it operates on individual
variance parameters directly (one-layer), not on a hierarchical decomposition.
There is no ``meta-uncertainty'' parameter analogous to $\kappa$.

Both are important references for the broader prior specification literature
but address different questions from the paper.


\section{A More Precise Account of Novelty}

The earlier sections of this note, written during the initial literature search,
overstated the novelty of the law of total variance as a prior specification
tool, and incorrectly positioned Manderson \& Goudie (2023) as addressing the
same problem. Successive revisions have sharpened the assessment. This section
presents the current understanding.

\subsection{The paper as a two-step pipeline}

The paper's methodology has two distinct steps. The novelty is concentrated
in the second.

\medskip
\noindent\textbf{Step 1: Quantity space $\to$ parameter space} (\S3.1). The
Jacobian correction $p_\theta(\theta) =
p_\vartheta(g(\theta))\sqrt{\det(J^\top J)}$ transforms a prior specified on
generated quantities into a prior on model parameters. This is a well-known
change-of-variables formula. Manderson \& Goudie (2023) solve the same problem
numerically, via distribution matching, for models where $g$ may not be
analytically tractable. \emph{Step 1 is not novel; the paper applies a known
technique and cites Segert (2019) for the derivation.}

\medskip
\noindent\textbf{Step 2: Direct prior $\to$ BHM hyperpriors} (\S3.2). This is
where the paper departs from existing work. Given a direct prior on $\vartheta$
for a single individual --- whether obtained via the Jacobian correction or
specified directly --- the researcher must now derive hyperpriors for the
population distribution. The integral to be satisfied is:
\[
  p^*(\vartheta) = \int p(\vartheta|\psi)\, p(\psi|\xi)\, d\psi
\]
This is an underdetermined problem: the direct prior provides two constraints
$(m^*, v^*)$, but the NIG hyperprior has four free parameters
$(\mu_0, \tau_0^2, \alpha_0, \beta_0)$. No existing work addresses this step.
Manderson \& Goudie stop after step~1. Fuglstad et al.\ (2020) decompose
variance across model components at a single level, not across hierarchical
levels. Kadane (1980) and Gabry et al.\ (2019) articulate the principle of
quantity-space specification but provide no mechanism for the BHM extension.

The paper resolves the underdetermination by introducing two interpretable
researcher choices:
\begin{itemize}
  \item $\rho$: what fraction of the direct prior variance reflects uncertainty
    about the population mean (vs.\ genuine heterogeneity).
  \item $\kappa$: how uncertain the researcher is about the level of
    heterogeneity itself.
\end{itemize}
Together, $(m^*, v^*, \rho, \kappa)$ provide four constraints for four unknowns,
and the solution is closed-form via the law of total variance. The individual
ingredients are known (the VPC/ICC ratio, the IG coefficient of variation, the
law of total variance), but their assembly into a determined system for
hyperprior derivation is new.

\subsection{What is known vs.\ what is new}

To be precise about each ingredient:

\begin{enumerate}
  \item \textbf{The law of total variance.} A textbook identity. Using it in the
    forward direction --- deriving the marginal moments of a NIG hierarchy ---
    is standard. Using it in the \emph{reverse} direction as a design equation
    (``given target moments, solve for hyperprior parameters'') is
    straightforward algebra but does not appear to have been written down as a
    methodology for prior specification. \emph{The technique is known; its
    application as a design tool for BHM priors is new.}

  \item \textbf{The $\rho$ parameterisation.} Numerically the same object as the
    VPC/ICC (Fisher 1925, Goldstein et al.\ 2002), which has been used
    exclusively as a posterior diagnostic. Fuglstad et al.\ (2020) use variance
    proportions as prior inputs, but in a different decomposition (across model
    components, with the proportions as random variables). Using the VPC/ICC as
    a \emph{fixed} design input for the between/within split in a BHM hierarchy
    is new. \emph{New role for a known quantity.}

  \item \textbf{The $\kappa$ parameterisation.} The CV of the IG distribution is
    textbook. Defining $\kappa = \SD[\sigma^2]/\E{\sigma^2}$ in a
    distribution-free way and proposing it as the primary specification device
    is new. \emph{New framing of a known relationship.}

  \item \textbf{The determined system.} The key observation is that $(\rho,
    \kappa)$ are not just interpretive glosses on existing parameters --- they
    are the two constraints needed to make the hyperprior derivation
    \emph{determined}. Without them, the researcher faces an underdetermined
    system (two moment equations in four unknowns) that can only be resolved by
    ad hoc choices of $\alpha_0$ and the variance split, or by numerical
    optimisation with regularisation. The paper replaces both with substantive
    meaning. \emph{This assembly is new.}
\end{enumerate}

\subsection{Manderson \& Goudie as the numerical Jacobian correction}

The initial drafts of this note contrasted Manderson \& Goudie (2023) with the
variance decomposition, framing them as ``numerical vs.\ analytical'' solutions
to the same problem. This was wrong --- twice. The first correction (Section
4.1.1) recognised that they operate at different levels of the hierarchy. The
second, sharper correction recognises that Manderson \& Goudie correspond to
the paper's \emph{first} step (the Jacobian correction), not the second
(the variance decomposition):

\begin{center}
\begin{tabular}{lll}
\toprule
\textbf{Pipeline step} & \textbf{This paper} & \textbf{Manderson \& Goudie} \\
\midrule
Beliefs about quantities & Jacobian correction & Distribution matching \\
$\;\to$ prior on $\theta$ & (analytical, \S3.1) & (numerical) \\
\addlinespace
Direct prior on $\vartheta$ & Variance decomposition & \emph{Not addressed} \\
$\;\to$ BHM hyperpriors $\xi$ & with $(\rho, \kappa)$ (\S3.2) & \\
\bottomrule
\end{tabular}
\end{center}

\noindent Both the Jacobian correction and Manderson \& Goudie solve the same
problem: ``I have beliefs about derived/observable quantities; translate these
into a prior on model parameters.'' The Jacobian correction does this
analytically when $g$ is injective and differentiable; Manderson \& Goudie do
it numerically for arbitrary models by matching the full prior predictive
distribution.

The paper's second step --- the variance decomposition with $(\rho, \kappa)$
--- has \emph{no counterpart} in Manderson \& Goudie. This is the step that
extends the methodology to BHMs, taking a single-individual direct prior and
deriving population-level hyperpriors. Without this step, a researcher with a
direct prior on $\vartheta$ would still face the $4p$ problem: how to specify
$(\mu_0, \tau_0^2, \alpha_0, \beta_0)$ for each parameter.

The two approaches \emph{could} be composed: Manderson \& Goudie (or the
Jacobian correction) for step~1, the variance decomposition for step~2. Neither
alone covers the full pipeline from quantity-space beliefs to BHM hyperpriors.

\subsection{What is the paper actually contributing?}

The paper makes two contributions of different character:

\begin{enumerate}[label=(\roman*)]
  \item \textbf{Step 1 (quantity space $\to$ parameter space)} is a known
    technique applied systematically to a new domain. The Jacobian correction is
    standard; Manderson \& Goudie (2023) provide the numerical analogue for
    non-conjugate models. The paper's contribution here is the careful
    application to structural models of risk preferences (proving injectivity
    of the CE mapping, handling the overdetermined case, etc.), not the
    technique itself. This step is best framed as \emph{applying} existing
    methodology.

  \item \textbf{Step 2 (direct prior $\to$ BHM hyperpriors)} is the
    methodological contribution. The problem it solves --- deriving hyperprior
    parameters from a single-individual direct prior --- is not addressed by
    any existing paper. The solution combines known ingredients (law of total
    variance, VPC/ICC, IG coefficient of variation) into a determined system
    via two interpretable parameters $(\rho, \kappa)$. This is more than a
    synthesis: it identifies a specific underdetermined problem in BHM prior
    specification and provides a principled, closed-form resolution. No
    existing approach does this.
\end{enumerate}

\noindent The gap that Step~2 fills is real and specific. Without it, a
researcher who has specified a direct prior on a generated quantity for a single
individual has no systematic way to derive the $4p$ hyperprior parameters needed
for a BHM. The available alternatives are:
\begin{itemize}
  \item \textbf{Ad hoc choices:} pick $\alpha_0$ and the variance split by
    convention or convenience (e.g.\ $\alpha_0 = 2.001$, equal split). This is
    common practice but provides no interpretive anchor.
  \item \textbf{Visual calibration:} simulate the prior predictive and adjust
    hyperparameters until the plots ``look right'' (Gabry et al., 2019). This
    is useful as a check but does not scale to $4p$ dimensions and provides no
    closed-form relationship.
  \item \textbf{Numerical optimisation:} extend Manderson \& Goudie's approach
    to the BHM setting. This is possible in principle but would face severe
    non-uniqueness (the hyperprior predictive is now two integrations removed
    from observables) and would require regularisation with no substantive
    interpretation.
\end{itemize}
The paper's $(\rho, \kappa)$ parameterisation replaces all three with a
determined, closed-form, interpretable derivation. Each degree of freedom in
the hyperprior corresponds to either a moment of the direct prior (which
encodes domain knowledge) or a structural belief about the hierarchy (which the
researcher can reason about independently of the parameter scale).


\section{Summary}

The paper's contribution is best understood as two steps of different character:

\begin{center}
\begin{tabular}{llll}
\toprule
\textbf{Step} & \textbf{What it does} & \textbf{Closest prior art} & \textbf{Novelty} \\
\midrule
Jacobian & Quantity space & Manderson \& Goudie & Known technique; \\
correction & $\to$ parameter & (2023): numerical & M\&G is the numerical \\
(\S3.1) & space prior & analogue & analogue \\
\addlinespace
Variance & Direct prior & \emph{None} & New: resolves the \\
decomposition & $\to$ BHM & & underdetermined \\
with $(\rho, \kappa)$ & hyperpriors & & $4p$ problem with \\
(\S3.2) & & & interpretable choices \\
\bottomrule
\end{tabular}
\end{center}

\noindent The individual ingredients of Step~2 have existing relatives:

\begin{center}
\begin{tabular}{lll}
\toprule
\textbf{Ingredient} & \textbf{Closest prior art} & \textbf{Nature of difference} \\
\midrule
$\rho$ (variance partition & VPC/ICC (Goldstein 2002); & Post-hoc diagnostic, not \\
\quad ratio) & HD priors (Fuglstad 2020) & design input; Fuglstad \\
 & & decomposes across model \\
 & & components, not hier.\ levels \\
\addlinespace
$\kappa$ (CV of & IG shape parameter & Known math; new as a named, \\
\quad heterogeneity) & (textbook) & distribution-free \\
 & & specification device \\
\addlinespace
Determined system: & \emph{None} & The assembly of $(\rho, \kappa)$ \\
$(m^*, v^*, \rho, \kappa)$ & & as the two constraints that \\
$\to (\mu_0, \tau_0^2,$ & & make the system determined \\
$\alpha_0, \beta_0)$ & & is new \\
\bottomrule
\end{tabular}
\end{center}


\section{Recommendations for the Paper}

\subsection{Framing the contribution}

The paper should be framed as having two contributions of different character:

\begin{enumerate}
  \item \textbf{Step 1 (quantity-space prior specification, \S3.1):} An
    application of the known Jacobian change-of-variables formula to structural
    models of risk preferences, with careful treatment of injectivity and the
    overdetermined case. This is best presented as \emph{applying} existing
    methodology to a new domain, not as a methodological innovation. Manderson
    \& Goudie (2023) should be cited as providing the numerical analogue for
    models where the analytical Jacobian is not available.

  \item \textbf{Step 2 (hyperprior derivation via variance decomposition,
    \S3.2):} The methodological contribution. The paper identifies a specific
    underdetermined problem in BHM prior specification --- the direct prior
    provides two moment constraints but the hyperprior has four free parameters
    --- and resolves it with two interpretable, unit-agnostic researcher
    choices $(\rho, \kappa)$ that make the system determined and yield
    closed-form hyperprior parameters. No existing paper addresses this
    problem. The individual ingredients (law of total variance, VPC/ICC, IG
    coefficient of variation) are known, but their assembly into a determined
    system for hyperprior derivation is new.
\end{enumerate}

\subsection{Citations to add}

\begin{enumerate}
  \item \textbf{Manderson \& Goudie (2023).} ``Translating Predictive
    Distributions into Informative Priors,'' arXiv:2303.08528. The numerical
    analogue of the paper's Step~1 (Jacobian correction). They translate
    beliefs about observables into prior parameters for single-level models via
    distribution matching and Bayesian optimisation. Cite in \S3.1 or \S2 as
    the numerical counterpart, noting that their approach applies to arbitrary
    models while the Jacobian correction requires an analytically tractable
    injective mapping. Note also that they face a non-uniqueness problem
    (resolved via generic regularisation), which motivates the paper's
    interpretable resolution via $(\rho, \kappa)$ at Step~2. The two
    approaches could be composed: M\&G for Step~1, the variance decomposition
    for Step~2.

  \item \textbf{Fuglstad, Hem, Knight, Rue \& Riebler (2020).} ``Intuitive
    Joint Priors for Variance Parameters,'' \emph{Bayesian Analysis}
    15(4):1109--1137. The closest existing work to Step~2, and the most
    important reference to cite and distinguish from. Likely reviewers may come
    from this group. Suggested framing: ``Fuglstad et al.\ (2020)
    reparameterise variance components into a total variance and proportions
    allocated across model components along a tree structure. The present
    approach differs in three respects: (i) it decomposes across hierarchical
    levels via the law of total variance rather than across model components;
    (ii) the proportion $\rho$ is a fixed researcher choice rather than a
    random variable with its own prior; and (iii) the derivation starts from a
    target unconditional distribution for a single individual.''

  \item \textbf{Goldstein, Browne \& Rasbash (2002).} ``Partitioning Variation
    in Multilevel Models,'' \emph{Understanding Statistics} 1(4):223--231. Cite
    when $\rho$ is introduced in \S3.2, noting the relationship to the VPC/ICC:
    same object, different role (design input vs.\ posterior diagnostic).

  \item \textbf{Aguilar \& B\"{u}rkner (2023).} ``The R2D2M2 Prior,''
    \emph{Electronic Journal of Statistics} 17(1):1319--1374. Cite in \S2 for
    the shared goal of making variance hyperparameters interpretable. Note the
    different domain (regression vs.\ hierarchical random effects) and the
    different decomposition (explained variance across predictors vs.\
    unconditional variance across hierarchical levels).

  \item \textbf{Natarajan \& Kass (2000).} ``Reference Bayesian Methods for
    Generalized Linear Mixed Models,'' \emph{JASA} 95:227--237. Cite for the
    induced-prior pathology: standard $\IG(\varepsilon, \varepsilon)$ priors on
    individual variance components induce bimodal, near-degenerate priors on
    the ICC. This strengthens the motivation for specifying $\rho$ directly
    rather than letting it emerge as an artefact of independent hyperprior
    choices.
\end{enumerate}

\subsection{Where to integrate}

\begin{enumerate}
  \item \textbf{Section 2 (Literature).} The lit review currently organises the
    prior specification literature into two camps. A third strand has emerged
    since those camps were defined: recent work using variance proportions or
    predictive distributions as inputs for prior specification (Fuglstad et al.\
    2020, Aguilar \& B\"{u}rkner 2023, Manderson \& Goudie 2023). A short
    paragraph positioning the paper relative to this strand would be valuable.
    The framing should be additive: these are contributions the paper builds
    alongside, not corrects. The paper's Step~1 is the analytical counterpart
    of Manderson \& Goudie; Step~2 fills a gap none of them address (the BHM
    extension).

  \item \textbf{Section 3.1 (Quantity-Space Prior Specification).} Cite
    Manderson \& Goudie as the numerical analogue of the Jacobian correction,
    noting that their approach handles non-conjugate models at the cost of
    requiring simulation and optimisation.

  \item \textbf{Section 3.2 (Two Interpretable Parameters).} When $\rho$ is
    introduced, a footnote noting the relationship to the VPC/ICC (Goldstein et
    al., 2002): same ratio, different role. When $\kappa$ is introduced, note
    that the underlying IG relationship is well known but the distribution-free
    framing and its use as a primary specification device are new. Cite
    Fuglstad et al.\ (2020) for the distinction: they decompose across model
    components with random proportions; the present paper decomposes across
    hierarchical levels with fixed choices.

  \item \textbf{Discussion (Section 8).} Frame the paper's position in the
    emerging toolkit:
    \begin{itemize}
      \item Manderson \& Goudie: beliefs about observables $\to$ prior
        parameters (numerical, single-level).
      \item Jacobian correction: beliefs about generated quantities $\to$ prior
        on parameters (analytical, requires injective $g$).
      \item Variance decomposition with $(\rho, \kappa)$: direct prior $\to$
        BHM hyperpriors (analytical, requires moment-matchable hierarchy).
      \item Fuglstad et al.: variance proportions across model components
        (random, with PC or Dirichlet priors).
    \end{itemize}
    These are complementary, not competing. The paper provides one piece ---
    the BHM extension --- that the others do not.
\end{enumerate}

\subsection{Terminology consideration}

The paper uses ``variance partition ratio'' for $\rho$. The established term is
``variance partition \emph{coefficient}'' (VPC). Two options:

\begin{enumerate}[label=(\alph*)]
  \item \textbf{Adopt ``coefficient.''} Gains instant familiarity for
    multilevel modelling readers. Risk: readers may assume $\rho$ is the same
    thing as the standard VPC (a posterior diagnostic), causing confusion about
    its role as a prior specification input.
  \item \textbf{Keep ``ratio.''} Signals that this is a distinct (though
    related) concept. Introduce it with a note: ``The variance partition ratio
    $\rho$ is numerically equivalent to the variance partition coefficient
    (Goldstein et al., 2002) but serves a different purpose here: it is a fixed
    researcher choice for prior construction, not a summary of a fitted model.''
\end{enumerate}

\noindent Option (b) is recommended. It avoids conflation while acknowledging
the connection.

\printbibliography

\end{document}
